Read in Gauge Theory Gravity, $f$ below is the metric function of the Schwarzschild-de Sitter field in the Painlev\'e-Gullstrand gauge, $\beta^{2} = 1 - f$, its Killing field is $K = \hb^{-1}(e_{t})$ and the curvature is the field strength (Fact \claimref{synthesis.imported-gtg-dictionary}, clauses (iv) and (vi)). \emph{Setting.} Units $\hbar = c = G = 1$ and $f(r) = 1 - 2M/r - \Lambda r^{2}/3$. For $\Lambda > 0$ and $9\Lambda M^{2} < 1$, $f$ has two positive zeros $r_{b} < r_{c}$ and a negative one $-(r_{b} + r_{c})$; by Vieta, $\Lambda/3 = 1/(r_{b}^{2} + r_{b}r_{c} + r_{c}^{2})$ and $2M = (\Lambda/3)\,r_{b}r_{c}(r_{b} + r_{c})$, and $f(r_{b}) = 0$ gives $\Lambda r_{b}^{2}/3 = 1 - u$ with $u = 2M/r_{b}$. Since $f(2M) = -4\Lambda M^{2}/3 < 0$ and $f$ is negative below $r_{b}$, $r_{b} > 2M$; and $f'(r_{b}) > 0$ gives $\kappa_{b} = M/r_{b}^{2} - \Lambda r_{b}/3 = (3M - r_{b})/r_{b}^{2} > 0$, so $r_{b} < 3M$ and $u \in (\tfrac23, 1)$. Inside $r_{b}$ near $r = 0$ the region is the interior of Theorem \claimref{interior.schwarzschild-de-sitter-weighted-limit}: same Kasner exponents, boundary data and wall as Schwarzschild, with corrections of relative order $\Lambda r^{3}/M$. So the placement rules (C1), (C2), and their leading-order form (C2$'$) on the exceptional set (Theorem \claimref{synthesis.exceptional-boundary-measure-one-calibration}), apply on this family as written.

\emph{(i)} Write $A(r) = -f(r) = (2M/r)\,U(r)$ with $U = 1 - r/2M + \Lambda r^{3}/6M$. The proper time to the wall is $\tau(r) = \int_{0}^{r}A^{-1/2}\,dr'$ and the three-volume per unit Killing time of the slice at $r$ is $V = 4\pi r^{2}\sqrt{A}$. With $r = \rho^{2}$, $\tau = (2M)^{-1/2}\int_{0}^{\rho}2\rho'^{2}U^{-1/2}\,d\rho'$ and $V = 4\pi(2M)^{1/2}\rho^{3}U^{1/2}$; both are power series in $\rho$ starting at $\rho^{3}$, with leading coefficients $\tfrac23(2M)^{-1/2}$ and $4\pi(2M)^{1/2}$, so $\mu_{E} = \lim V/\tau = 12\pi M$. $\Lambda$ appears in $U$ only through $\rho^{6} = r^{3}$, hence in $V$ and $\tau$ only from relative order $\Lambda r^{3}/M$, and not in $\mu_{E}$. On the slice $r_{c}^{3} = 2M$, $(4\pi r_{c}^{2})^{2}A(r_{c}) = 16\pi^{2}r_{c}^{4}(2M/r_{c})(1 - r_{c}/2M + \Lambda r_{c}^{3}/6M) = 64\pi^{2}M^{2}(1 - r_{c}/2M + \Lambda/3)$, using $r_{c}^{3} = 2M$, which is the stated rate; Proposition \claimref{singularity.stretched-singularity-count} is its $\Lambda = 0$ case. In Planck units the corrections are $O((2M)^{-2/3})$ and $O(\Lambda)$; both vanish in the limit $M \to \infty$, $\Lambda \to 0$ at fixed $\Lambda M^{2}$, in which $u$ stays fixed.

\emph{(ii)} With the epoch $\Delta t$ and $\nu'$ modes per Planck area of $\mu_{E}$, $N = \nu'\mu_{E}\Delta t = 12\pi\nu' M\,\Delta t$, exactly for (C2$'$); for (C1) and (C2), $N = 8\pi\nu M\,\Delta t$ at leading order, the same number with $\nu = \tfrac32\nu'$ (Theorem \claimref{synthesis.exceptional-boundary-measure-one-calibration}(iii)). The products below do not depend on the normalisation of $\partial_{t}$, since $\mu_{E}$ is per unit $t$ and the periods are in $t$. At $\Lambda = 0$, Theorem \claimref{scrambling.epoch-modular-period-two-unknowns-become} gives the epoch $8\pi M$ under (IC) and the area law fixes the density: $\nu'\langle s\rangle_{0}\cdot 12\pi M\cdot 8\pi M = 4\pi M^{2}$, so $\nu'\langle s\rangle_{0} = 1/24\pi$ (in the 1-nat convention, $\langle s\rangle_{0} = 1$, $\nu' = 1/24\pi$; in the capacity convention $\langle s\rangle_{0} = \ln 2$). With the hole's period $\Delta t = 2\pi/\kappa_{b}$,
\[
\frac{N\langle s\rangle_{0}}{\pi r_{b}^{2}} = \frac{1}{24\pi}\cdot\frac{12\pi M\cdot 2\pi}{\kappa_{b}\,\pi r_{b}^{2}} = \frac{M}{\kappa_{b}r_{b}^{2}} = \frac{M}{3M - r_{b}} = \frac{u}{3u - 2} .
\]
With $w_{h} = 2M/r_{b}^{3}$, the radial eigenvalue of the electric part of the Weyl tensor, $2\Lambda/(3w_{h}) = \Lambda r_{b}^{3}/3M = (1 - u)\,r_{b}/M = 2(1 - u)/u$, so $1 - 2\Lambda/3w_{h} = (3u - 2)/u$ and the ratio is $1/(1 - 2\Lambda/3w_{h})$. It equals $1$ at $u = 1$ ($\Lambda = 0$); its derivative in $u$ is $-2/(3u - 2)^{2} < 0$, and $u$ decreases as $\Lambda M^{2}$ grows ($r_{b}$ increases with $\Lambda$ at fixed $M$, since $\partial_{\Lambda}f < 0$ and $f'(r_{b}) > 0$), so the ratio increases, and it diverges as $u \to \tfrac23$, $r_{b} \to 3M$, $9\Lambda M^{2} \to 1$, where $\kappa_{b} \to 0$. For small $\Lambda M^{2}$, $u = 1 - \tfrac43\Lambda M^{2} + O(\Lambda^{2}M^{4})$ and the ratio is $1 + \tfrac83\Lambda M^{2} + O(\Lambda^{2}M^{4})$. Dividing the count's entropy by $A_{b} = 4\pi r_{b}^{2}$ gives $\eta_{\mathrm{count}} = \tfrac14/(1 - 2\Lambda/3w_{h})$; $\Lambda = \tfrac14 R$ is the Ricci part and $w_{h}$ the Weyl eigenvalue of the curvature at every point of the horizon.

\emph{(iii)} $N\langle s\rangle_{0} = \pi r_{b}^{2}$ requires $12\pi M\Delta t/24\pi = \pi r_{b}^{2}$, that is $\Delta t = 2\pi r_{b}^{2}/M$. Since $r_{b} > 2M$, $\kappa_{b} = M/r_{b}^{2} - \Lambda r_{b}/3 < M/r_{b}^{2}$. For the cosmological horizon, scale $r_{b} = 1$ and write $x = r_{c} > 1$: then $\Lambda/3 = 1/(1 + x + x^{2})$, $M = x(1 + x)/(2(1 + x + x^{2}))$ and $\kappa_{c} = \Lambda r_{c}/3 - M/r_{c}^{2} = (2x^{2} - x - 1)/(2x(1 + x + x^{2}))$, so $\kappa_{c} < M$ is $2x^{2} - x - 1 < x^{2}(1 + x)$, that is $x^{3} - x^{2} + x + 1 > 0$, true for $x \ge 1$ since $x^{3} - x^{2} \ge 0$. With $\kappa_{c} < \kappa_{b}$ (the check verifies it on five rational families with both horizons rational; it is not needed for the conclusion, which uses only $\kappa_{c}, \kappa_{b} < M/r_{b}^{2}$), both periods exceed $2\pi r_{b}^{2}/M$ and both counts exceed $A_{b}/4$.

\emph{(iv)} Under one mode set, the hole's horizon modes are its cutoff modes, so the count's entropy $N\langle s\rangle$ is the hole's horizon entropy. Theorem \claimref{synthesis.sds-clausius-fixes-area-constant}(i),(ii) gives $S_{b} = A_{b}/4$ under the series' temperatures, the Clausius relation along the family and $S_{b} \to 0$ as $M \to 0$. With $\langle s\rangle = \langle s\rangle_{0}$ the two disagree by the factor of (ii) for every $\Lambda > 0$. If instead the counted modes are required to carry $A_{b}/4$ (the area-law hypothesis of Theorem \claimref{synthesis.residue-from-mode-entropy}), then $N\langle s\rangle = A_{b}/4$ with $N$ as in (ii) gives $\langle s\rangle = \langle s\rangle_{0}(3u - 2)/u$. By Theorem \claimref{synthesis.boundary-mode-thermal-entropy}, $s(x)$ is even, strictly decreasing in $|x|$ and at most $\ln 2$, with $x = 2\pi\omega/\kappa_{b}$; a smaller mean entropy needs larger $|x|$ for some counted modes. So the weights in units of $\kappa_{b}$ cannot be the same on all members of the family, and if they are all zero at $\Lambda = 0$ (the registered value, the equality case of Theorem \claimref{synthesis.residue-from-mode-entropy}(ii)) they are not all zero for any $\Lambda > 0$. A law $\omega_{k} = \kappa_{b}k$ fixes every $x_{k} = 2\pi k$ and hence $\langle s\rangle$, so it cannot follow $(3u - 2)/u$. The same holds for every weight law the store states: the zero weight of the registered residue (Theorem \claimref{synthesis.residue-from-mode-entropy}(ii)), $\kappa k$, $\kappa(k + \tfrac12)$ and the values $0, \pm\kappa$ of the unadopted clock hypothesis are all fixed multiples of the surface gravity, each fixing $\langle s\rangle$ on the family (by a search of the store's weight assignments; none involves $\Lambda$ except through $\kappa$). Under Premise (Q), whose clause (Q-w) sets the slope $c = \kappa$, no assignment escapes. The affine class $\omega_{k} = ck + d$ of Theorem \claimref{synthesis.q-splits-maximum-entropy-affine-weights}(ii) leaves $c$ free, so without (Q-w) a slope depending on $\Lambda/w_{h}$ could follow $(3u - 2)/u$; no stored result adopts one, and on modes shared by the two sheets of an edge (R$_{\mathrm{loc}}$) forces zero weight (filed alongside, synthesis.r-loc-needs-sheet-labelled-modes).

\emph{(v)} By (ii) the count, read as an entropy per unit area of the hole's horizon, depends on $\Lambda/w_{h}$, so it is not one constant (on a one-parameter family this curvature form is a relabelling of the dependence on $u$; what matters is that the ratio varies), and in particular not a density depending only on the flat tangent data. A derivation of (Loc) from (C1), (C2) and (IC), read with a horizon's modular period, with one mode set would have to make the count a fixed multiple of the area on every family to which these rules apply; Schwarzschild-de Sitter is such a family and the multiple varies. A curvature-independent density is obtained by placing modes on the horizon with a density fixed by flat data (per unit area and unit rapidity), which is (Loc) restated in the language of modes. (IC) is stated in the store for Schwarzschild (epoch $8\pi M$); reading it on this family with a horizon's modular period is the assumption of (ii), and (iii) shows that no horizon period removes the disagreement. As for the literal Kerr count (Theorem \claimref{synthesis.kerr-ring-count-misses-area-law}, cited as a parallel and not used), the conflict is between the cutoff count and $A_{b}/4$, and it bears on the literal cutoff placement and on (IC) with a horizon period away from $\Lambda = 0$, not on (Loc), whose value $\tfrac14$ agrees with $A_{b}/4$.

\emph{Remarks.} (a) The physical size is small: for a hole of the mass of the Sun in Planck units ($M \sim 10^{38}$) and the observed $\Lambda \sim 10^{-122}$, $\tfrac83\Lambda M^{2} \sim 10^{-46}$; the result is structural, not observational. (b) The statement uses the cosmological constant only as a Ricci curvature at the horizon; it does not use any premise about the frozen point. (c) The cosmological horizon carries its own modes in the store (with an unfixed count), and is not part of this comparison.

\emph{The check.} Exact power series (GA.ps\_pow\_rat, GA.ps\_int) for $\tau$ and $V$ in $\rho$ give $\mu_{E} = 12\pi M$ for four masses with $2M$ a square and four values of $\Lambda$ (including $0$ and a negative one), and show that $\Lambda$ first moves the series at relative order $\rho^{6}$; the (C1) volume identity is checked at three rational $r_{c}$ and three $\Lambda$. The ratio of (ii) is computed from $\nu' = 1/24\pi$, $\mu_{E}$ and $2\pi/\kappa_{b}$, with $f(r_{b}) = 0$ and $\kappa_{b} = f'(r_{b})/2$ checked, on five rational $r_{b}$ from $2M$ towards $3M$ (values $1, \tfrac{10}{9}, 2, 10, 100$), and compared with $u/(3u - 2)$ and $1/(1 - 2\Lambda/3w_{h})$; towards Nariai it exceeds $10^{j}/3$ at $r_{b} = 3M - M/10^{j}$. On five rational families with both horizons rational the needed epoch is $2\pi r_{b}^{2}/M$ and $0 < \kappa_{c} < \kappa_{b} < M/r_{b}^{2}$; the polynomial inequality of (iii) is checked on $192$ ratios. The forced mean entropy $(3u - 2)/u$ is checked on four members.
